%==============================================================================
% FICHAMENTO CRÍTICO — Effectiveness of Creating Digital Twins with Different
% Digital Dentition Models and CBCT
% Conforme NBR 6023 (referências) e NBR 10520 (citações)
%==============================================================================
\section{Fichamento: \citeonline{lee2023twins}}

% --- 1. IDENTIFICAÇÃO ---
\subsection{Identificação}
\begin{tabular}{ll}
\textbf{Título:} & Effectiveness of creating digital twins with different
digital dentition models and cone-beam computed tomography \\
\textbf{Autor(es):} & Lee, Joo-Hee; Lee, Hye-Lim; Park, In Young; On, Sung-Woon;
Byun, Soo-Hwan; Yang, Byoung-Eun \\
\textbf{Periódico:} & Scientific Reports \\
\textbf{Ano:} & 2023 \\
\textbf{Volume:} & 13 \\
\textbf{Número:} & 1 \\
\textbf{Páginas:} & 10603 \\
\textbf{DOI:} & \url{https://doi.org/10.1038/s41598-023-37774-x} \\
\textbf{Qualis:} & A1 (Multidisciplinar) \\
\textbf{Tipo:} & Artigo original \\
\end{tabular}

% --- 2. OBJETIVO ---
\subsection{Objetivo}
Avaliar a eficácia de diferentes modelos de dentição digital (escaneamento de
moldes de alginato em diferentes intervalos de tempo e escaneamento intraoral
com diferentes métodos) como alternativas ao modelo de gesso para criar gêmeos
digitais odontológicos quando combinados com CBCT, quantificando as discrepâncias
tridimensionais em pontos de referência específicos.

% --- 3. METODOLOGIA ---
\subsection{Metodologia}
\textbf{Desenho do estudo:} Estudo transversal comparativo intra-sujeito com 5
grupos de gêmeos digitais por paciente. \\
\textbf{Amostra:} 20 pacientes (9 meninos, 11 meninas, 12--18 anos) com dentição
permanente completa, sem síndromes craniofaciais ou aparelhos ortodônticos
metálicos. \\
\textbf{Métricas principais:} Diferença média (em mm) entre grupos em 6 pontos
de referência; voxel size do CBCT (0,39 mm) como limiar de aceitação clínica. \\
\textbf{Validação:} Análise de sobreposição 3D via programa R2GATE\texttrademark;
comparação entre 5 métodos por paciente (controle: gesso + CBCT; Grupos I--IV:
alginato 5 min, alginato 2 h, IOS segmentado CS 3600, IOS arco completo i700).

% --- 4. RESULTADOS PRINCIPAIS ---
\subsection{Resultados Principais}
\begin{itemize}
  \item Todos os grupos apresentaram diferenças médias inferiores ao voxel size
        do CBCT (0,39 mm), indicando acurácia clinicamente aceitável.
  \item O Grupo II (alginato após 2 horas) apresentou as maiores discrepâncias
        ($>$ 0,2 mm em todos os pontos de referência).
  \item O Grupo I (alginato em até 5 min) e o Grupo III (IOS segmentado CS 3600)
        apresentaram as menores diferenças (máximo de 0,053 mm).
  \item A varredura de alginato dentro de 5 minutos mostrou-se equivalente ao
        IOS segmentado, ambos superiores ao alginato após 2 horas.
\end{itemize}

% --- 5. LIMITAÇÕES ---
\subsection{Limitações}
\begin{itemize}
  \item Amostra pequena (n = 20) e restrita a adolescentes sem anomalias
        craniofaciais, limitando a generalização para outras faixas etárias.
  \item O estudo utilizou apenas a arcada maxilar; não foram avaliados os
        arcos mandibulares nem a oclusão dinâmica.
  \item As condições de armazenamento do alginato (temperatura, umidade) foram
        controladas em ambiente laboratorial, não refletindo necessariamente
        as condições clínicas reais.
\end{itemize}

% --- 6. CONTRIBUIÇÃO PARA O LIVRO ---
\subsection{Contribuição para o Livro}
Este artigo é a referência central do Capítulo 21, que aborda gêmeos digitais
em odontologia. O estudo fornece evidências quantitativas sobre a acurácia de
diferentes métodos de aquisição de modelos dentários para compor gêmeos
digitais, demonstrando que alternativas digitais (IOS e escaneamento de
alginato) são clinicamente equivalentes ao modelo de gesso tradicional.
O conceito de voxel size do CBCT como limiar de aceitação clínica (0,39 mm)
é utilizado no livro para estabelecer critérios de qualidade para a fusão de
imagens na criação de gêmeos digitais. A comparação entre diferentes métodos
de escaneamento intraoral (segmentado vs. arco completo) orienta a escolha do
protocolo de aquisição mais adequado.

% --- 7. ANÁLISE CRÍTICA ---
\subsection{Análise Crítica}
\begin{itemize}
  \item \textbf{Validade interna:} Alta -- Desenho intra-sujeito com 5 grupos
        por paciente, controle das condições de armazenamento e uso de
        software comercial de sobreposição.
  \item \textbf{Validade externa:} Baixa -- Amostra pequena (n = 20), faixa
        etária restrita (12--18 anos), apenas maxila, sem patologias ou
        reabilitações protéticas.
  \item \textbf{Reprodutibilidade:} Alta -- Protocolo detalhado com marcas
        comerciais específicas (CS 3600, i700, Alphard 3030, R2GATE),
        parâmetros de aquisição e pós-processamento claramente documentados.
  \item \textbf{Relevância clínica:} Alta -- Orienta a escolha entre métodos
        de aquisição digital para planejamento cirúrgico virtual, confecção
        de guias cirúrgicos e wafers oclusais.
  \item \textbf{Conflito de interesses:} Não -- Os autores declararam ausência de
        conflitos.
  \item \textbf{Viés identificado:} Viés de seleção -- Amostra de conveniência
        com pacientes sem restaurações metálicas ou aparelhos ortodônticos,
        condição que não reflete a maioria dos pacientes que necessitam de
        gêmeos digitais.
\end{itemize}

% --- 8. CITAÇÕES RELEVANTES ---
\subsection{Citações Relevantes}
\begin{citacao}
``The mean differences in all groups were less than the 0.39 mm CBCT voxel
size. Although there may be slight differences in accuracy between groups,
all were clinically acceptable.'' (p. 8).
\end{citacao}

% --- 9. MAPEAMENTO SDD ---
\subsection{Mapeamento SDD}
\textbf{Problema:} Gêmeos digitais odontológicos dependem de fusão precisa entre
CBCT e modelos dentários; o modelo de gesso tem limitações logísticas e
alternativas digitais precisam ser validadas. \\
\textbf{Entrada:} Imagens CBCT (DICOM, voxel 0,39 mm) + modelos dentários
digitais (STL) de 5 origens: gesso, alginato 5 min, alginato 2 h, IOS
segmentado CS 3600, IOS arco completo i700. \\
\textbf{Saída:} Gêmeo digital fusionado com diferenças quantificadas em 6
pontos de referência por grupo, em milímetros. \\
\textbf{Critérios:} Diferença média $<$ voxel size do CBCT (0,39 mm);
discrepância máxima $\leq$ 0,2 mm para métodos alternativos viáveis;
equivalência estatística entre método alternativo e controle (gesso). \\
\textbf{Confiança:} 0,72 -- Estudo meticuloso, mas com amostra pequena e
restrita que limita a generalização dos achados para a prática clínica
heterogênea.

% --- 10. REFERÊNCIA ABNT ---
\subsection{Referência ABNT}
\begin{verbatim}
LEE, Joo-Hee et al. Effectiveness of creating digital twins with different
digital dentition models and cone-beam computed tomography. Scientific
Reports, London, v. 13, n. 1, p. 10603, 2023. DOI:
10.1038/s41598-023-37774-x.
\end{verbatim}
